\section{Results}
\label{sec:results}

\subsection{RQ1 --- Endpoint coverage}

The LLM suite exercises \RQIopsCovered{} of \RQIopsTotal{} operations:
\textbf{\RQIcoveragePct\%} coverage, uniform across all three subjects
(\texttt{rest-ncs} $\RQINcsCovered/\RQINcsTotal$, \texttt{rest-scs}
$\RQIScsCovered/\RQIScsTotal$, \texttt{features-service}
$\RQIFeaturesCovered/\RQIFeaturesTotal$). The one-sample Wilcoxon signed-rank test against
the 90\% threshold gives $W=\RQIwilcoxonW$, $p=\RQIpvalue$, rank-biserial
$=\RQIrankBiserial$. $H_0$ is \textbf{\RQIreject} at $\alpha=\alphaRaw$, and also under the
stricter Bonferroni threshold $\alphaBonf$.

The pre-registered miss profile by endpoint type is empty: with every operation covered, no
endpoint type is missed. This is a real result but a weak one --- it says the LLM can
\emph{reach} every documented operation, not that it tests any of them well. RQ2 is where
that distinction bites.

\subsection{RQ2 --- Fault detection against ground truth}

This is the pre-registered primary comparison, and it does not go the LLM's way.

Against the \RQIIn{} seeded faults, overall Recall is \textbf{LLM \RQIIrecallLlm},
\textbf{manual \RQIIrecallManual}, and \textbf{EvoMaster \RQIIrecallEvo} --- in absolute
terms, \PDkilledLlm{}, \PDkilledManual{}, and \PDkilledEvo{} mutants killed out of
\PDkilledTotal{}. EvoMaster detects twice as many faults as the LLM.

The Friedman test across the three arms over per-subject Recall gives
$\chi^2=\RQIIfriedmanChi$, $p=\RQIIfriedmanP$, so the pre-registered omnibus test does not
reject. With $n=3$ subjects this test is severely underpowered, which is why the proposal
designated the pooled per-fault McNemar test as the primary RQ2 analysis.

That test compares the arms fault-by-fault over all \RQIIn{} mutants. For LLM versus
EvoMaster the discordant pairs are $b=\RQIImcnemarLEb$ (killed by the LLM only) against
$c=\RQIImcnemarLEc$ (killed by EvoMaster only): $\chi^2=\RQIImcnemarLEstat$,
$p=\RQIImcnemarLEp$, Cliff's $\delta=\RQIIcliffLE$. $H_0$ is \textbf{\RQIImcnemarLEreject}.
The LLM is not significantly \emph{worse} than EvoMaster, but the point estimate is
unambiguously in EvoMaster's favour and the effect size is negative. Approximate achieved
power is \RQIImcnemarLEpower{}, so this is an underpowered non-result rather than evidence
of equivalence --- we can rule out neither a real gap nor its absence, and we say so rather
than reading $p>\alpha$ as parity.

\begin{table}[t]
\centering
\caption{RQ2 --- fault-detection Recall per subject and overall. EvoMaster leads overall;
LLM and manual are identical everywhere (Section~\ref{sec:results-degenerate}).}
\label{tab:rq2}
\small
\begin{tabular}{@{}lrrr@{}}
\toprule
\textbf{Subject} & \textbf{LLM} & \textbf{Manual} & \textbf{EvoMaster} \\
\midrule
\texttt{rest-ncs} (70 mutants)        & \RQIINcsLlm          & \RQIINcsManual      & \textbf{\RQIINcsEvo} \\
\texttt{rest-scs} (59 mutants)        & \textbf{\RQIIScsLlm} & \RQIIScsManual      & \RQIIScsEvo \\
\texttt{features-service} (4 mutants) & \RQIIFeaturesLlm     & \RQIIFeaturesManual & \RQIIFeaturesEvo \\
\midrule
\textbf{Overall} (\RQIIn{} mutants)   & \RQIIrecallLlm       & \RQIIrecallManual   & \textbf{\RQIIrecallEvo} \\
\bottomrule
\end{tabular}
\end{table}

Table~\ref{tab:rq2} shows the overall ordering is not uniform. EvoMaster's lead is carried
entirely by \texttt{rest-ncs} (\RQIINcsEvo{} against \RQIINcsLlm{}), the most
arithmetic-heavy subject; on \texttt{rest-scs} the LLM leads (\RQIIScsLlm{} against
\RQIIScsEvo{}), and on \texttt{features-service} all three arms tie.
Section~\ref{sec:discussion} argues this is exactly what the oracle-kind account predicts.

\subsubsection{A degenerate comparison: LLM versus manual}
\label{sec:results-degenerate}

The LLM-versus-manual contrast must be reported, but it cannot be interpreted. The two arms
kill \emph{identical} mutant sets on every subject: the McNemar table has
$b=\RQIImcnemarLMb$ and $c=\RQIImcnemarLMc$ --- \emph{zero} discordant pairs --- giving
$p=\RQIImcnemarLMp$ and Cliff's $\delta=\RQIIcliffLM$ by construction rather than by
measurement. The Holm-adjusted pairwise Wilcoxon likewise returns $p=\RQIIholmLMp$ on
identical per-subject vectors.

We attribute this to the two suites sharing an oracle \emph{kind}
(Section~\ref{sec:discussion}) and, critically, to a construct limitation we do not want to
understate: in this run both suites were authored by the same agent under different role
instructions. Two suites drawn from one author, asserting on the same spec-derived status
codes, are not independent observers of the fault population. The RQ2 conclusion that
survives is therefore \emph{LLM versus EvoMaster}; the LLM-versus-manual leg is reported for
completeness and excluded from interpretation. Section~\ref{sec:threats} states what would
be needed to repair it.

\subsection{RQ3 --- Edge-case scenarios per endpoint}

The LLM produces \textbf{\RQIIItotalLlm{}} edge-case scenarios against manual's
\textbf{\RQIIItotalManual{}} across \RQIIIn{} operations (median \RQIIImedianLlm{} versus
\RQIIImedianManual{} per operation). The paired one-tailed Wilcoxon signed-rank test gives
$W=\RQIIIwilcoxonW$, $p=\RQIIIpvalue$, rank-biserial $=\RQIIIrankBiserial$. $H_0$ is
\textbf{\RQIIIreject}, and survives the Bonferroni threshold $\alphaBonf$.

The effect is consistent rather than driven by outliers: the LLM produces more scenarios on
\RQIIIpairsLlm{} of \RQIIIn{} operations, manual leads on \RQIIIpairsManual{}, and
\RQIIIpairsTie{} tie. This is the largest effect in the study. It is also, as
Section~\ref{sec:results-pvd} shows, the least consequential.

\subsection{Produced versus detected}
\label{sec:results-pvd}

Table~\ref{tab:pvd} places the volume and effectiveness metrics side by side, and the two
orderings disagree.

\begin{table}[t]
\centering
\caption{What each arm \emph{produced} versus what it \emph{detected}. The volume ordering
(LLM $>$ manual $>$ EvoMaster) inverts on the ground truth (EvoMaster $>$ LLM $=$ manual).}
\label{tab:pvd}
\small
\begin{tabular}{@{}lrrr@{}}
\toprule
 & \textbf{LLM} & \textbf{Manual} & \textbf{EvoMaster} \\
\midrule
\multicolumn{4}{@{}l}{\emph{Produced}} \\
Test cases                        & \textbf{\PDtestsLlm}     & \PDtestsManual        & \PDtestsEvo \\
Edge-case scenarios               & \textbf{\RQIIItotalLlm}  & \RQIIItotalManual     & --- \\
\midrule
\multicolumn{4}{@{}l}{\emph{Detected (against ground truth)}} \\
Mutants killed / \PDkilledTotal   & \PDkilledLlm             & \PDkilledManual       & \textbf{\PDkilledEvo} \\
HTTP error behaviours / \PDerrKey & \textbf{\PDerrTriggeredLlm} & \PDerrTriggeredManual & \PDerrTriggeredEvo \\
5xx responses elicited            & \textbf{\PDcrashLlm}     & \PDcrashManual        & \PDcrashEvo \\
\bottomrule
\end{tabular}
\end{table}

The LLM writes \PDtestsLlm{} test cases to EvoMaster's \PDtestsEvo{} ---
$\PDvolumeRatioLlmEvo\times$ the volume --- and kills half as many mutants
(EvoMaster detects $\PDkillRatioEvoLlm\times$ the LLM's faults). It does lead on the
error-surface dimension,
triggering \PDerrTriggeredLlm{} of the \PDerrKey{} documented HTTP error behaviours against
manual's \PDerrTriggeredManual{} and EvoMaster's \PDerrTriggeredEvo{}, and eliciting
\PDcrashLlm{} 5xx responses including a genuine server crash. So the LLM is not weak
everywhere: it is strong at reaching the \emph{documented} error surface and weak at
detecting \emph{value} faults.

Had we reported only the metrics the surveyed literature uses --- coverage, scenario counts,
500-error tallies --- every one of them would have ranked the LLM first, and the conclusion
would have been that it beats both baselines. The ground truth reverses that. This is the
practical case for GAP-D: without a denominator, the field's headline comparisons are
measuring volume and reporting it as effectiveness.

\begin{figure}[t]
  \centering
  \includegraphics[width=\columnwidth]{fig1_distribution.png}
  \caption{RQ2 --- distribution of fault-detection Recall across the three arms over
  \RQIIn{} seeded faults. EvoMaster's advantage is concentrated on the arithmetic-heavy
  \texttt{rest-ncs} subject.}
  \label{fig:rq2}
\end{figure}

\begin{figure}[t]
  \centering
  \includegraphics[width=\columnwidth]{fig2_comparison.png}
  \caption{RQ3 --- per-endpoint edge-case scenario counts, LLM versus manual across
  \RQIIIn{} operations. The LLM leads on \RQIIIpairsLlm{} operations, manual on
  \RQIIIpairsManual{}.}
  \label{fig:rq3}
\end{figure}
